\section{Setting and design problem}\label{sec:setting}

\subsection{The model in brief}

The MPME sequence~\cite{cheng2019} is an unbalanced steady-state gradient echo that records
three configuration pathways ($k=+1,0,-1$), three echoes each, in each of two scans. The two
scans share the repetition time $\TR=25$\,ms and the readout and differ only in the nominal flip
angle; the actual angle of scan $i$ is $x_i=\Bone\alpha_i$. The unknowns are
$\theta=(\ln\Mzero,\ln\Bone,\ln\Tone,\ln\Ttwo,\ln\Ttwos,\dw,\phi_0)$, with logarithms for the positive
parameters so that standard deviations are relative. We use the following results of the
companion technical note~\cite{technote}.

\begin{itemize}
\item \emph{Two invariants per scan.} Each pathway amplitude is $a_k=v\,G_k(u,E_2,\dw)$, where
  $u$ and $v$ depend on the flip angle and $\Tone$ only, and $G_k$ does not depend on either.
\item \emph{Ernst-angle form.} With the Ernst angle $\alpha_E=\arccos\E^{-\TR/\Tone}$ and
  $\zeta=\tan(\alpha_E/2)=\sqrt{\tanh(\TR/2\Tone)}$, define the \emph{offset} of a sample from the Ernst angle
  \begin{equation}
    w=\ln|\tan(x/2)|-\ln\zeta .
    \label{eq:w}
  \end{equation}
  Then $u=\tanh w$ and $|v|=\zeta\operatorname{sech}w$, with $\sgn v=\sgn\tan(x/2)$. On the axis
  $s=\ln|\tan(x/2)|$ the spoiled-gradient-echo amplitude is a $\operatorname{sech}$ bump centred at the
  Ernst angle; a scan samples it at $s$.
\item \emph{What the data determine.} After the echo decay and phases have fixed $\Ttwo$, $\Ttwos$,
  $\dw$ and $\phi_0$, one scan determines its offset $w$ and the amplitude $\ell=\ln(\Mzero\zeta)$; two
  scans determine $w_1$, $w_2$ and the shared $\ell$, three numbers for $(\Mzero,\Bone,\Tone)$. One scan
  alone cannot determine $\Bone$.
\end{itemize}

\subsection{The design problem}\label{sec:problem}

\paragraph{Design variables.} The nominal angles $0<\alpha_1<\alpha_2\le540^\circ$; $\TR$ and the
readout are fixed, so every design has the same acquisition time.

\paragraph{Criterion.} The Cram\'er--Rao bound (CRLB) on $\ln\Tone$ with all seven parameters free,
in units of $\sigma/\Mzero$ (complex Gaussian noise, standard deviation $\sigma$ per real and
imaginary part). Bounds on $\ln\Bone$, $\ln\Ttwo$ and $\ln\Mzero$ are reported alongside.

\paragraph{Robustness range.} The design must work for every $\Bone\in[0.7,1.3]$ (the brain at
3\,T, central to peripheral) and the five tissues of \cref{tab:tissues}. We report the
\emph{worst case} of the bound over this range and its \emph{mean} (uniform over $\Bone$, equal
weight per tissue).

\begin{table}[!htbp]
\centering
\caption{Tissues of the robustness range ($\Mzero=1$ for all, so bounds are relative to each
tissue's own signal). Ernst angles at $\TR=25$\,ms.}
\label{tab:tissues}
\small
\begin{tabular}{lcccc}
\toprule
Tissue & $\Tone$ (ms) & $\Ttwo$ (ms) & $\Ttwos$ (ms) & $\alpha_E$ ($^\circ$)\\
\midrule
CSF & 4000 & 1500 & 900 & 6.4\\
grey matter (GM) & 1350 & 90 & 60 & 11.0\\
white matter (WM) & 850 & 66 & 50 & 13.8\\
deep GM & 1130 & 58 & 40 & 12.0\\
lesion & 1500 & 120 & 80 & 10.4\\
\bottomrule
\end{tabular}
\end{table}

\paragraph{Pipelines.} Two uses of the data lead to different optima.
\begin{itemize}
\item \emph{Joint fit}: both scans at full resolution, all parameters fitted per voxel; the
  criterion is the joint $\Tone$ bound above.
\item \emph{Two-stage}: a short, low-resolution scan 2 gives a smooth $\Bone$ map whose noise is
  negligible after smoothing~\cite[Sec.~5]{technote}, and $\Tone$ comes from scan 1 with $\Bone$
  known. The criterion for $\alpha_1$ is the scan-1 $\Tone$ bound with $\Bone$ known; $\alpha_2$ only needs to
  determine $\Bone$ reliably.
\end{itemize}

\paragraph{Constraints.} The scans are acquired one after the other, so the specific absorption
rate (SAR) of each scan scales as the square of its own angle for a fixed pulse, and a
peak-$\Bone$ limit fixes the shortest duration of the largest pulse. Both constrain the largest
angle $\alpha_{\max}=\alpha_2$; we report the best design as a function of $\alpha_{\max}$.

\paragraph{Identifiability.} A design must not admit a complex alias (\cref{def:alias}) for any true
$\Bone$ in the robustness range within the fitting range $\Bone\in[0.6,1.4]$; this is checked analytically
(\cref{sec:aliases}). Magnitude aliases are allowed: complex data, the FID-sign test or, for some
designs, a narrower fitting range remove them.

\paragraph{Model error.} Finally, the shortlisted designs are compared for their bias under
errors the fitting model ignores: a pulse-ratio error, finite pulse duration, slice or slab
profile (\cref{sec:robustness}).
